% Generated draft: requires mathematical and source review.
\documentclass[11pt]{article}
\usepackage[margin=25mm]{geometry}
\usepackage{fontspec}
\setmainfont{Latin Modern Roman}
\usepackage{amsmath,amssymb,mathtools,unicode-math}
\setmathfont{Latin Modern Math}
\usepackage{enumitem,needspace,xurl,hyperref}
\setlength{\parindent}{0pt}
\setlength{\parskip}{5pt}
\setlength{\emergencystretch}{3em}
\newcommand{\N}{\mathbb{N}}
\newcommand{\Z}{\mathbb{Z}}
\newcommand{\Q}{\mathbb{Q}}
\newcommand{\R}{\mathbb{R}}
\newcommand{\C}{\mathbb{C}}
\newcommand{\F}{\mathbb{F}}
\newcommand{\A}{\mathbb{A}}
\newcommand{\PP}{\mathbb{P}}
\newcommand{\E}{\mathbb{E}}
\newcommand{\eps}{\varepsilon}
\newcommand{\dd}{\,\mathrm{d}}
\begin{document}
\section*{UnsolvedMath ID: 30003160}
\textbf{Source problem code:} catalogue entry
\subsection{Definitions and mathematical statement}
In the usual one-space-dimensional interpretation, a \(2\times2\) system of conservation laws has the form
\[
\partial_t u+\partial_x f(u)=0,\qquad u(0,x)=u_0(x),
\]
where \(u\) has two real components and \(f\) is a flux taking values in \(\R^2\). The source does not specify the state space, spatial domain, time interval, flux regularity, structural hypotheses, boundary conditions, or admissible class of solutions. This equation therefore records the conventional setting suggested by the source, rather than a fully specified initial or initial--boundary value problem.

A family \(\{\Gamma_\alpha\}_{\alpha\in A}\) of curves in a space--time domain \(D\) is locally finite if every compact subset of \(D\) intersects only finitely many of the curves. For a discontinuity along a differentiable graph \(x=s(t)\), with left and right traces \(u^-\) and \(u^+\), the conservation-law jump condition is
\[
 s'(t)\bigl(u^+-u^-\bigr)=f(u^+)-f(u^-).
\]
The admissibility conditions distinguishing the intended shocks are not supplied.

\textbf{Source question, with unresolved qualifications.} For \(2\times2\) systems of conservation laws, including isentropic gas dynamics, do generic \(C^3\) initial data give rise to solutions that are smooth outside a locally finite family of shock curves? Can one prove uniqueness and continuous dependence of the resulting general solutions?

Here \(C^3\) means three times continuously differentiable on the initial spatial domain. The source does not specify the admissible data space or its topology, the meaning of ``generic,'' or the differentiability required by ``smooth.'' In particular, a residual-set interpretation of genericity and a \(C^\infty\) interpretation of smoothness are not asserted. The domain \(D\) on which local finiteness is required, including its relation to the initial time and any spatial boundary, is also unspecified. Uniqueness requires a specified solution and admissibility class; continuous dependence requires specified topologies for data and solutions and a time range. These choices cannot be recovered from the supplied statement. No pressure law, density restrictions, or treatment of vacuum is specified for the isentropic gas example. Thus the record supplies a research question whose precise hypotheses and quantifiers remain incomplete.
\subsection{Short English statement}
The question asks whether generic \(C^3\) initial data for two-component conservation laws, including isentropic gas dynamics, produce solutions smooth away from locally finitely many shock curves, and whether the resulting general solutions are unique and depend continuously on the data. The system class, genericity, solution admissibility, regularity, domains, and dependence topologies require clarification before this becomes a fully quantified conjecture.
\subsection{Sources}
\begin{itemize}
\item Supplied local UnsolvedMath registry record, permanent ID \(30003160\), title ``Generic Regularity of Two-by-Two Conservation Laws,'' exact source problem code \texttt{catalogue entry}; collection: Oberwolfach Reports Open Problems. No individual proposer is identified; the supplied proposed year is 2016.
\item Supplied source citation: \emph{Hyperbolic Techniques in Modelling, Analysis and Numerics} (2017), Oberwolfach Reports, pp.\ 1683--1751, DOI \texttt{10.4171/owr/2016/30}. \url{https://doi.org/10.4171/owr/2016/30}. The background attributes the question to workshop open-problem discussions documented in Oberwolfach Reports. The supplied citation says 2017, whereas the DOI and proposed-year metadata refer to 2016; this discrepancy is retained without resolution.
\item Supplied literature source: ``Oberwolfach Report 30.'' \url{https://ems.press/content/serial-article-files/46634?nt=1}.
\item Supplied literature source: ``Chen, generic regularity of conservative solutions.'' \url{https://gengchen.ku.edu/sites/gengchen/files/documents/genericwave22.pdf}. No further bibliographic details or precise applicability to this question are supplied.
\item Literature assessment checked on \texttt{2026-08-21}: the supplied curation labels the problem open and states that generic regularity and uniqueness or continuous dependence for general two-by-two systems, including the isentropic gas system, remain outstanding beyond established special cases. This is supplied dated curation, not an independently verified literature assessment; the special cases are not identified. The background explicitly requests independent verification of time-sensitive status and cited claims.
\item The registry labels the statement \texttt{corrected\_verified}, explaining that only notation, typography, and TeX-delimiter cleanup was applied and that hypotheses, quantifiers, and subquestions were preserved. This supplied label does not resolve the formulation gaps or establish mathematical correctness.
\end{itemize}
\end{document}
